\chapter{Known Cases and Evidence}\label{known-cases}
\chapterstatus{complete}

\section{Introduction}\label{known-cases:section-introduction}

This chapter collects the cases in which the Hodge conjecture is known, and explains the mechanism in each. They fall into three groups. First, cases where every Hodge class
is forced to be a combination of known algebraic classes: small dimension, cellular varieties, varieties whose Hodge classes are generated by divisors, and hypersurfaces outside the middle
degree. Second, abelian varieties, where the Mumford--Tate group (Chapter~\ref{mumford-tate}) makes the Hodge classes computable; there the difficulty concentrates in the \emph{Weil classes}.
Third, geometric methods: the decomposition of the diagonal for varieties with small Chow groups, normal functions for cubic fourfolds, and the Kuga--Satake construction for K3 surfaces.
We state precisely what is proved, with proofs of the elementary parts and references for the rest. References: the source \emph{The Hodge Conjecture} (Chapter 6),
\cite{lewis-survey}, \cite{birkenhake-lange}, \cite{voisin-hodge-2}.

\noindent\textbf{Prerequisites.} Chapter~\ref{lefschetz-11} (The Lefschetz (1,1) Theorem), Chapter~\ref{mumford-tate} (Mumford--Tate Groups), Chapter~\ref{vhs} (Variations of Hodge Structure) and
Chapter~\ref{intermediate-jacobians} (Intermediate Jacobians and Normal Functions).

Throughout, $X$ is a nonsingular complex projective variety of dimension $n$.

\section{Cases from general theory}\label{known-cases:section-general}

\begin{definition}\label{known-cases:definition-lefschetz-classes}
The \emph{Lefschetz ring} of $X$ is the $\QQ$-subalgebra $\mathrm{Lef}^\bullet(X) \subseteq \bigoplus_pH^{2p}(X; \QQ)$ generated by $\mathrm{Hdg}^1(X)$, the rational classes of divisors. Its elements are
\emph{Lefschetz classes}.
\end{definition}

\begin{proposition}\label{known-cases:proposition-lefschetz-classes}
Lefschetz classes are algebraic. Hence $\mathrm{HC}(X)$ holds whenever $\mathrm{Hdg}^\bullet(X) = \mathrm{Lef}^\bullet(X)$.
\end{proposition}

\begin{proof}
Divisor classes are algebraic by the Lefschetz $(1,1)$-theorem (Theorem~\ref{lefschetz-11:theorem-lefschetz-11}), and products of algebraic classes are algebraic
(Theorem~\ref{cycle-class:theorem-cycle-class}(1)).
\end{proof}

\begin{theorem}\label{known-cases:theorem-general}
The Hodge conjecture holds in the following cases:
\begin{enumerate}
\item $p \in \{0, 1, n - 1, n\}$, and hence for $\dim X \leq 3$ (Corollary~\ref{lefschetz-11:corollary-codimension-n-1});
\item cellular varieties, such as flag varieties and toric varieties (Example~\ref{hodge-conjecture:example-cellular});
\item very general abelian varieties and all their powers (Theorem~\ref{mumford-tate:theorem-very-general});
\item smooth complete intersections, in all degrees except the middle one (Proposition~\ref{known-cases:proposition-hypersurfaces});
\item uniruled fourfolds \cite{conte-murre-1978} and, more generally, varieties whose group of $0$-cycles is supported on a subvariety of dimension $\leq 3$, in codimension $2$
\cite{bloch-srinivas-1983};
\item cubic fourfolds \cite{zucker-1977};
\item Fermat hypersurfaces of many degrees and dimensions \cite{shioda-1979}.
\end{enumerate}
\end{theorem}

Items (1)--(4) are proved in this book. For (5)--(7) see the references and Sections~\ref{known-cases:section-diagonal} and~\ref{known-cases:section-hypersurfaces}.

\section{Abelian varieties}\label{known-cases:section-abelian}

\begin{remark}\label{known-cases:remark-abelian-mechanism}
For an abelian variety $A$, the Hodge classes on all powers $A^r$ are the invariants of the Hodge group $\mathrm{Hg} = \mathrm{Hg}(H^1(A))$ (Corollary~\ref{mumford-tate:corollary-hodge-ring-abelian}),
and the divisor classes on $A^r$ are the invariants of the \emph{Lefschetz group} $\mathrm{Lef}(A)$, the connected centraliser of $\End(A) \otimes \QQ$ in the symplectic group of a polarisation
\cite{murty-1984}. Always $\mathrm{Hg} \subseteq \mathrm{Lef}(A)$. If equality holds, the Hodge ring of every power of $A$ is generated by divisors, and the Hodge conjecture holds for all powers of $A$.
\end{remark}

\begin{theorem}\label{known-cases:theorem-abelian-known}
$\mathrm{Hg}(A) = \mathrm{Lef}(A)$, and hence the Hodge conjecture holds for all powers of $A$, in each of the following cases: $A$ is a product of elliptic curves \cite{imai-1976},
\cite{murty-1984}; $A$ is simple of prime dimension (Tankeev) \cite{tankeev-1982}; $A$ is of dimension at most $5$ and has no factor of Weil type (Moonen--Zarhin) \cite{moonen-zarhin-1999}. More
generally the Hodge conjecture holds for $A$ itself (not necessarily its powers) whenever $\mathrm{Hdg}^\bullet(A) = \mathrm{Lef}^\bullet(A)$, which holds for $A$ of dimension $\leq 3$ by
Theorem~\ref{known-cases:theorem-general}(1).
\end{theorem}

\begin{proposition}[Weil classes]\label{known-cases:proposition-weil-classes}
Let $A$ be a complex abelian fourfold and $K \subseteq \End(A) \otimes \QQ$ an imaginary quadratic field such that, for the two embeddings $\sigma, \bar\sigma \colon K \to \CC$, the eigenspaces of $K$
on $H^{1,0}(A)$ both have dimension $2$ (\emph{Weil type}). Then $H^1(A; \QQ)$ is a $4$-dimensional $K$-vector space. Let $V_\sigma, V_{\bar\sigma} \subseteq H^1(A; \CC)$ be the eigenspaces of $K$.
There is a $2$-dimensional $\QQ$-subspace $W_K \subseteq \Lambda^4H^1(A; \QQ) = H^4(A; \QQ)$ with $W_K \otimes \CC = \Lambda^4V_\sigma \oplus \Lambda^4V_{\bar\sigma}$, and it consists of Hodge classes, the
\emph{Weil classes}. For a general abelian fourfold of Weil type, $\mathrm{Hdg}^2(A) = \mathrm{Lef}^2(A) \oplus W_K$ with $W_K \cap \mathrm{Lef}^2(A) = 0$, so the Weil classes are
not combinations of intersections of divisors.
\end{proposition}

\begin{proof}
$K$ acts on $H_\CC = H^1(A; \CC)$ with eigenspaces $V_\sigma$, $V_{\bar\sigma}$, each of dimension $4$, and $K$ acts by Hodge endomorphisms, so $V_\sigma = V_\sigma^{1,0} \oplus V_\sigma^{0,1}$. By
hypothesis $\dim V_\sigma^{1,0} = 2$, and then $\dim V_\sigma^{0,1} = 2$. The space $\Lambda^4_KH^1(A; \QQ) \otimes_\QQ\CC$ decomposes as $\Lambda^4V_\sigma \oplus \Lambda^4V_{\bar\sigma}$, and
$\Lambda^4V_\sigma = \Lambda^2V_\sigma^{1,0} \otimes \Lambda^2V_\sigma^{0,1}$ is of type $(2,2)$, and likewise for $\bar\sigma$. The subspace $\Lambda^4V_\sigma \oplus \Lambda^4V_{\bar\sigma}$ of
$\Lambda^4H^1(A; \CC)$ is defined over $K$, since the eigenspace decomposition is, and it is stable under the nontrivial automorphism of $K$, which exchanges $V_\sigma$ and $V_{\bar\sigma}$. By Galois
descent it is the complexification of a $\QQ$-subspace $W_K$, of dimension $2$, whose elements are rational classes of type $(2,2)$. That $W_K \cap \mathrm{Lef}^2 = 0$ for general $A$ of Weil type, where $\End(A) \otimes \QQ = K$ and the Lefschetz
group is a unitary group $\mathrm{U}(2,2)$ while $\mathrm{Hg} = \mathrm{SU}(2,2)$, is \cite{weil-1977}; see also \cite{birkenhake-lange}.
\end{proof}

\begin{remark}\label{known-cases:remark-weil-classes}
Weil classes are the essential difficulty for abelian varieties: they are the first Hodge classes not generated by divisors, and Moonen and Zarhin showed that in dimension $\leq 5$
all exceptional Hodge classes come from Weil classes. The Hodge conjecture for Weil classes was proved by Schoen for $K = \QQ(\sqrt{-3})$ and $K = \QQ(i)$ in special cases \cite{schoen-1988},
by van Geemen in further cases, and by Markman for all abelian fourfolds of Weil type of discriminant $1$ \cite{markman-2023}. Weil classes are absolute Hodge classes by Deligne's theorem
(Chapter~\ref{absolute-hodge}).
\end{remark}

\section{Hypersurfaces and complete intersections}\label{known-cases:section-hypersurfaces}

\begin{proposition}\label{known-cases:proposition-hypersurfaces}
Let $X \subseteq \PP^{n+1}$ be a nonsingular hypersurface, or more generally a nonsingular complete intersection of dimension $n$ in some $\PP^N$, and $h$ the hyperplane class. Then for $2p \neq n$,
$H^{2p}(X; \QQ) = \QQ h^p$. Hence $\mathrm{HC}(X, p)$ holds for $2p \neq n$.
\end{proposition}

\begin{proof}
By the Lefschetz hyperplane theorem (Theorem~\ref{kodaira:theorem-lefschetz-hyperplane}), applied repeatedly for complete intersections, $H^k(X; \QQ) \cong H^k(\PP^N; \QQ)$ for $k < n$, which is
$\QQ h^{k/2}$ for $k$ even. For $k > n$, Poincar\'e duality gives $b_k = b_{2n-k}$, which is $1$ for $k$ even, and $h^{k/2} \neq 0$ by Proposition~\ref{kahler:proposition-kahler-class-powers}, so
$H^k = \QQ h^{k/2}$. Powers of $h$ are algebraic, being classes of complete intersections of hyperplane sections.
\end{proof}

\begin{remark}\label{known-cases:remark-middle-degree}
In the middle degree $n = 2p$, a very general hypersurface of degree $d$ has $\mathrm{Hdg}^p(X) = \QQ h^p$ as soon as $H^{n}$ is not of Tate type up to $h^p$, for instance for surfaces with $d \geq 4$
(Theorem~\ref{lefschetz-11:theorem-noether-lefschetz}) and for fourfolds with $d \geq 3$, so the conjecture holds there trivially. Special hypersurfaces are the interesting cases. For cubic
fourfolds, Zucker proved the conjecture with normal functions (Remark~\ref{intermediate-jacobians:remark-poincare-lefschetz}), and Voisin proved the integral version \cite{voisin-2006}. For
Fermat hypersurfaces $x_0^d + \cdots + x_{n+1}^d = 0$, the Hodge classes are computable via the action of $(\mu_d)^{n+2}$ and characters (Jacobi sums), and Shioda proved the conjecture for many $(d, n)$,
including $d$ prime and $d = 4$ \cite{shioda-1979}.
\end{remark}

\section{Uniruled varieties and the decomposition of the diagonal}\label{known-cases:section-diagonal}

\begin{theorem}[Bloch--Srinivas]\label{known-cases:theorem-bloch-srinivas}
Suppose that $\CH_0(X)_\QQ$ is supported on a closed subvariety of dimension $\leq 3$, that is, there is $W \subseteq X$ of dimension $\leq 3$ with $\CH_0(W)_\QQ \to \CH_0(X)_\QQ$ surjective. Then
$\mathrm{HC}(X, 2)$ holds. This applies to rationally connected varieties, for which $\CH_0(X)_\QQ = \QQ$.
\end{theorem}

This is \cite{bloch-srinivas-1983}. The method is the \emph{decomposition of the diagonal}: the hypothesis gives an equality $N\Delta_X = Z_1 + Z_2$ in $\CH^n(X \times X)$ with $Z_1$ supported on
$D \times X$ for a divisor $D$ and $Z_2$ supported on $X \times W$. Letting both sides act on $H^4(X; \QQ)$ as correspondences, the identity (times $N$) factors through the cohomology of a resolution of
$D$, where Hodge classes of degree $2$ on a variety of lower dimension are handled by induction and the $(1,1)$-theorem, plus a contribution through $W$, of dimension $\leq 3$, where the conjecture is
known. The same argument gives the conjecture for uniruled fourfolds, first proved by Conte and Murre \cite{conte-murre-1978}.

\section{K3 surfaces, hyper-K\"ahler varieties and Kuga--Satake}\label{known-cases:section-k3}

\begin{remark}\label{known-cases:remark-k3}
For a single K3 surface the conjecture holds, by dimension. The interesting questions are on products. \emph{Mukai's conjecture}, proved by Buskin \cite{buskin-2019}, says that every rational Hodge
isometry $H^2(S; \QQ) \to H^2(S'; \QQ)$ between K3 surfaces, a Hodge class in $H^4(S \times S'; \QQ)$, is algebraic. Markman extended this to hyper-K\"ahler varieties of $K3^{[n]}$-type.
\end{remark}

\begin{remark}[Kuga--Satake]\label{known-cases:remark-kuga-satake}
Kuga and Satake attached to a polarised K3 surface $S$ an abelian variety $\mathrm{KS}(S)$, of dimension $2^{19}$, built from the even Clifford algebra of the lattice of primitive classes, together with
an embedding of Hodge structures $H^2_{\mathrm{prim}}(S; \QQ) \hookrightarrow H^1(\mathrm{KS}(S)) \otimes H^1(\mathrm{KS}(S))$ \cite{kuga-satake-1967}. The Hodge conjecture predicts that this embedding is
induced by an algebraic correspondence between $S$ and $\mathrm{KS}(S)^2$ (the \emph{Kuga--Satake Hodge conjecture}). It is known for some families, for instance by Paranjape for double planes branched
along six lines \cite{paranjape-1988}. Deligne used the Kuga--Satake construction, which works in families and in characteristic $p$, to prove the Weil conjectures for K3 surfaces
\cite{deligne-k3}, and to show that Hodge classes on K3 surfaces and their products with abelian varieties behave like absolute Hodge classes.
\end{remark}

\section{The evidence assessed}\label{known-cases:section-assessment}

\begin{remark}\label{known-cases:remark-assessment}
The known cases share a feature: the Hodge classes are either generated by divisors, or produced by an explicit geometric construction adapted to the variety (normal functions for cubic fourfolds,
the decomposition of the diagonal, families of cycles on special abelian fourfolds, derived equivalences of K3 surfaces). There is no general mechanism. At the same time no case has failed, the
Hodge loci behave as the conjecture predicts (Theorem~\ref{vhs:theorem-cdk}), and Hodge classes on abelian varieties satisfy the arithmetic consequence of the conjecture (Deligne,
Chapter~\ref{absolute-hodge}). The first genuinely open cases include Weil classes on general abelian fourfolds of Weil type of arbitrary discriminant, abelian varieties of dimension $\geq 6$ with
exceptional classes, and products of two K3 surfaces in degree $4$ away from Hodge isometries.
\end{remark}

\section{Exercises}\label{known-cases:section-exercises}

\begin{exercise}\label{known-cases:exercise-cubic-fourfold-numbers}
For a nonsingular cubic fourfold $X \subseteq \PP^5$, show that $H^4(X; \QQ)$ has Hodge numbers $(0, 1, 21, 1, 0)$ and that $\mathrm{HC}(X, p)$ holds for $p \neq 2$ by
Proposition~\ref{known-cases:proposition-hypersurfaces}.
\end{exercise}

\begin{solution}
As in Proposition~\ref{calabi-yau:proposition-quintic}, $c(TX) = (1 + h)^6/(1 + 3h)$, whose degree-$4$ part is $(15 - 3 \cdot 20 + 9 \cdot 15 - 27 \cdot 6 + 81)h^4 = 9h^4$, and
$\int_Xh^4 = 3$, so $\chi(X) = 27$. The odd Betti numbers vanish and $b_0 = b_2 = b_6 = b_8 = 1$ (Proposition~\ref{known-cases:proposition-hypersurfaces}), so $b_4 = 27 - 4 = 23$. The Hodge
numbers $h^{4,0} = 0$, $h^{3,1} = 1$, $h^{2,2} = 21$ follow from Griffiths' description of the primitive cohomology of hypersurfaces by residues \cite{voisin-hodge-2}. The second statement is
Proposition~\ref{known-cases:proposition-hypersurfaces} with $n = 4$.
\end{solution}

\begin{exercise}\label{known-cases:exercise-lefschetz-ring-product}
Let $E$ be an elliptic curve without complex multiplication. Show that $\mathrm{Hdg}^\bullet(E^r) = \mathrm{Lef}^\bullet(E^r)$ for all $r$.
\end{exercise}

\begin{solution}
By Proposition~\ref{mumford-tate:proposition-elliptic-curves}, $\mathrm{Hg}(E) = \mathrm{SL}_2$, and the invariants of $\mathrm{SL}_2$ acting diagonally on $\Lambda^\bullet(H^{\oplus r})$ are generated, by classical
invariant theory, by the invariants in $\Lambda^2(H^{\oplus r})$, which are the divisor classes (Exercise~\ref{abelian-varieties:exercise-hodge-classes-torus}). Apply
Corollary~\ref{mumford-tate:corollary-hodge-ring-abelian}.
\end{solution}
